\documentclass[11pt]{article}
\usepackage[utf8]{inputenc}
\usepackage{amsmath,amssymb}
\usepackage{hyperref}
\title{Robust Earthquake Early Warning Machine Learning under Distribution Shift}
\author{Anonymous}
\date{}
\begin{document}
\maketitle

\begin{abstract}
This paper studies Earthquake Early Warning Machine Learning. Grounded in a real, citable literature \cite{ref1,ref2,ref3,ref4,ref5,ref6,ref7,ref8},
we propose a focused method and report \textbf{illustrative} experiments.
All references below are real and verifiable (OpenAlex/arXiv); the reported
numbers are simulated to illustrate intended behaviour.
\end{abstract}

\section{Introduction}
Earthquake Early Warning Machine Learning is an active research area. This work targets a concrete gap identified from
the recent literature and contributes a method together with an evaluation protocol.

\section{Related Work (real literature)}
The following works, retrieved from public bibliographic APIs, frame our study:
\begin{itemize}
\item Kong et al. (2018) — "Machine Learning in Seismology: Turning Data into Insights", Seismological Research Letters (cited 500).
\item Li et al. (2018) — "Machine Learning Seismic Wave Discrimination: Application to Earthquake Early Warning", Geophysical Research Letters (cited 305).
\item Mousavi et al. (2019) — "A Machine‐Learning Approach for Earthquake Magnitude Estimation", Geophysical Research Letters (cited 297).
\item Khan et al. (2023) — "A Systematic Review of Disaster Management Systems: Approaches, Challenges, and Future Directions", Land (cited 190).
\item Rafiei et al. (2017) — "NEEWS: A novel earthquake early warning model using neural dynamic classification and neural dynamic optimization", Soil Dynamics and Earthquake Engineering (cited 171).
\item Lomax et al. (2019) — "An Investigation of Rapid Earthquake Characterization Using Single‐Station Waveforms and a Convolutional Neural Network", Seismological Research Letters (cited 158).
\end{itemize}

\section{Method}
We model distribution shift explicitly and optimise a worst-case objective over an uncertainty set of plausible test distributions. We instantiate this idea for Earthquake Early Warning Machine Learning and analyse its cost/quality trade-off.

\section{Experiments (Illustrative)}
\begin{center}
\begin{tabular}{lcc}
System & Primary Metric & Efficiency \\ \hline
Strong baseline & 0.812 & 1.00$\times$ \\
Ablation & 0.828 & 0.92$\times$ \\
\textbf{Ours} & \textbf{0.851} & \textbf{0.71$\times$}
\end{tabular}
\end{center}
\emph{Metrics simulated to illustrate the intended effect; not measured.}

\section{Conclusion}
We connected a real literature to a concrete method for Earthquake Early Warning Machine Learning. Future work will
validate the illustrative results with real experiments.

\begin{thebibliography}{9}
\bibitem{ref1} Qingkai Kong, Daniel T. Trugman, Zachary E. Ross et al.. Machine Learning in Seismology: Turning Data into Insights. \emph{Seismological Research Letters}, 2018. DOI:10.1785/0220180259
\bibitem{ref2} Zefeng Li, Men‐Andrin Meier, Egill Hauksson et al.. Machine Learning Seismic Wave Discrimination: Application to Earthquake Early Warning. \emph{Geophysical Research Letters}, 2018. DOI:10.1029/2018gl077870
\bibitem{ref3} S. Mostafa Mousavi, Gregory C. Beroza. A Machine‐Learning Approach for Earthquake Magnitude Estimation. \emph{Geophysical Research Letters}, 2019. DOI:10.1029/2019gl085976
\bibitem{ref4} Saad Mazhar Khan, Imran Shafi, Wasi Haider Butt et al.. A Systematic Review of Disaster Management Systems: Approaches, Challenges, and Future Directions. \emph{Land}, 2023. DOI:10.3390/land12081514
\bibitem{ref5} Mohammad Hossein Rafiei, Hojjat Adeli. NEEWS: A novel earthquake early warning model using neural dynamic classification and neural dynamic optimization. \emph{Soil Dynamics and Earthquake Engineering}, 2017. DOI:10.1016/j.soildyn.2017.05.013
\bibitem{ref6} Anthony Lomax, Alberto Michelini, Dario Jozinović. An Investigation of Rapid Earthquake Characterization Using Single‐Station Waveforms and a Convolutional Neural Network. \emph{Seismological Research Letters}, 2019. DOI:10.1785/0220180311
\bibitem{ref7} Martijn van den Ende, Jean‐Paul Ampuero. Automated Seismic Source Characterization Using Deep Graph Neural Networks. \emph{Geophysical Research Letters}, 2020. DOI:10.1029/2020gl088690
\bibitem{ref8} Men‐Andrin Meier, Zachary E. Ross, Anshul Ramachandran et al.. Reliable Real‐Time Seismic Signal/Noise Discrimination With Machine Learning. \emph{Journal of Geophysical Research Solid Earth}, 2018. DOI:10.1029/2018jb016661
\end{thebibliography}
\end{document}
